\section{Decisions and undecidability in nilpotent groups}\label{sec:nilpotent}

\subsection{Direct decomposition with torsion (N5)}\label{sec:n5}
\entry{N5}{G.~Baumslag}
Is direct indecomposability decidable for finitely generated nilpotent
groups? The torsion-free case is prior~\cite{bmo}. Here the input is a
finite presentation \emph{promised to define a nilpotent group}; no
recognition algorithm for arbitrary presentations is claimed.

\begin{theorem}\label{prop:n5}
Under this promise there is a uniform algorithm deciding whether
$G=A\times B$ with both factors nontrivial, and constructing generators
for such factors when they exist.
\end{theorem}
\begin{proof}
Use standard effective nilpotent-group operations: consistent
polycyclic presentations, subgroup membership, centres, finite torsion
subgroups, kernels and intersections. Let $T=T(G)$, $Z=Z(G)$,
$Q=G/T$, and let $R$ be the rational Malcev completion of $Q$.
Rational direct decomposition is effective. Uniqueness modulo the
centre of indecomposable rational factors~\cite[Theorem 3.3]{fgh},
also used in~\cite{bmo}, says that any decomposition matches a partition
of a fixed decomposition $R=R_1\times\cdots\times R_m$, modulo $Z(R)$.

Put $K=G/Z$ and let $\kappa:K\to R/Z(R)$ be induced by the natural
map. Its kernel is $S=T(K)$. To see this, the preimage $C$ in $G$
of $Z(Q)=Q\cap Z(R)$ satisfies $\ker\kappa=C/Z$. Choose a finite
quotient injective on $T$, using residual finiteness, and let $e$ be
its exponent. The characteristic power subgroup $G^e$ is torsion-free
and of finite index. For $c\in C$, $[c^e,g]$ lies in both $T$ and
$G^e$ for every $g$, so $c^e\in Z$. Conversely a central power has
central image in torsion-free $Q$, by unique roots. Thus $C/Z=T(K)$,
which is finite and computable.

For each ordered partition $I,J$ of the rational factors compute
\[
 K_1=\kappa^{-1}(R_I Z(R)/Z(R)),\qquad
 K_2=\kappa^{-1}(R_J Z(R)/Z(R)).
\]
If $G=A\times B$ induces this partition, its images $Y_1,Y_2$ in
$K$ form a direct product and satisfy $K_i=Y_iS$. Hence
$[K_i:Y_i]\le|S|$. Enumerate all subgroups of $K_i$ up to that
index, using finite permutation representations, and retain exactly
the pairs with
\[
 K=Y_1Y_2,\qquad[Y_1,Y_2]=1,\qquad Y_1\cap Y_2=1.
\]
This is a finite complete list of quotient decompositions that might
lift. Empty partitions and trivial $Y_i$ are kept, so central and
finite factors are not lost.

Let $X_i$ be the full preimages in $G$. Discard a pair if
$[X_1,X_2]\ne1$. Otherwise $X_1X_2=G$ and $X_1\cap X_2=Z$.
Choose presentations of $Y_i$ and lift their generators to $X_i$.
The lifted relators are central defects $v_i\in Z^{m_i}$; changing
lifts by central elements $z$ changes them to $v_i+E_i z$, where
$E_i$ is the exponent-sum matrix. For a proposed splitting
$Z=Z_1\oplus Z_2$, the extension $X_i/Z_i\to Y_i$ has a section
exactly when $E_i z=-v_i$ is solvable over $Z/Z_i$.

Smith reduction converts this into finitely many conditions
\begin{equation}\label{eq:central-conditions}
 v\in Z_i+dZ,\qquad d\ge0,
\end{equation}
where $d=0$ means exact membership. Appendix~\ref{app:central}
gives a finite constructive decision for these conditions over an
arbitrary finitely generated abelian group, retaining torsion, integral
splittings and required nontrivial factors. Corrected lifts and $Z_i$
then generate factors $A_i$. They commute, intersect trivially and
generate $G$. Every actual decomposition occurred in the finite
quotient list and passes precisely these lifting conditions, proving
completeness and termination.
\end{proof}

The rational decomposition machinery is credited prior work; the
proposed additional scope is torsion. In particular, rational
decomposability alone does not decide integral decomposability. The
final implementation accepts promised-nilpotent finite presentations
without a supplied class bound and returns original-generator factor
words. Its negative integral-gluing controls and finite comparisons
support implementation details, not the universal support argument.
\evidence{N5}{proof.md}

\subsection{Single commutators in every free nilpotent group (N8)}\label{sec:n8}
\entry{N8}{A.~Miasnikov}
Part (a) concerns $[x,y]=g$ in general finitely generated class-two
groups and has a prior negative answer~\cite{romankov2016}. Part (b) asks the same question
in every finitely generated free nilpotent group. We use
$[x,y]=x^{-1}y^{-1}xy$ throughout this section.

\begin{theorem}\label{prop:n8}
Given finite rank $r$, class $c$ and a word $g$ in
$N_{r,c}=F_r/\gamma_{c+1}(F_r)$, the equation $[x,y]=g$ is
decidable; on a positive answer its factors can be constructed.
\end{theorem}

The main issue is the finite treatment of correction parameters, not
enumeration of positive witnesses. Here is the proposed algorithm's
structure, with its separation argument detailed in
Appendix~\ref{app:commutator}.

For $g\ne1$, let $d$ be its first nonzero lower-central degree and
$W$ its leading integral Lie coordinate. Exact Nielsen operations on
the pair preserve its commutator. They normalize every solution to
leading terms
\[
 C\in(L_\Z)_p,\quad D\in(L_\Z)_q,\quad [C,D]=W,
 \quad p\le q,\quad p+q=d.
\]
Here $L_\Z$ is the untruncated free integral Lie ring. A finite list
of leading pairs suffices. For unequal weights, the projective bracket
map has finite fibres; rational points and all signed integral scale
divisors are enumerated. For equal weights, the exterior tensor gives
a rank-two lattice, and finitely many Hermite representatives suffice
modulo the exact $\SL_2(\Z)$ operations. Primitive normalization alone
would discard valid solutions.

At offset $s$ the new coordinate map is
\begin{equation}\label{eq:n8-correction}
 A_s(U,V)=[U,D]+[C,V].
\end{equation}
The structural kernel analysis separates two cases. After the offset
$q-p$, its elements lie in $\Lie(C,D)$ and have exact full-rank
integral period lattices, obtained by denominator clearing in rational
commutator-preserving substitutions. Every residue is retained. Before
that offset a nonzero kernel has dimension at most one; these are the
exceptional parameters.

At the first exceptional offset $t$, solve \emph{all} equations from
$t$ through $2t-1$ as one affine integer system. If the remaining class
ends before $2t$, all remaining equations are jointly linear and the
branch terminates by integer linear algebra. Otherwise choose an
integral parameter $T$ for the first exceptional coordinate in the
complete surviving lattice. Its image can have a nonunit step, which
must be retained. At offset $2t$ there is a scalar quadratic obstruction
in $T$ whose quadratic coefficient survives after quotienting by both
the new correction image and \emph{all later parameter columns}.
It has at most two integral roots. Fix the prefix through $t$ at
each root and restart, leaving higher coordinates free again.

The diagonal positional-polynomial lemma in
Appendix~\ref{app:commutator} is the reason later parameters cannot
cancel this quadratic term. Each branch has finitely many children
and every restart fixes a higher offset, so the recursion terminates.
An assembled solution is checked in the actual group. Completeness
follows by retaining every leading branch, every exact-period residue
and every integral root forced by the full correction block.

The classical inputs include Hall/Malcev coordinates, Shirshov--Witt,
Lazard elimination and the two-coefficient inner-solution theorem for
free Lie algebras~\cite{altassan,remeslennikov-stoehr}. They are applied
in suitable free subalgebras \emph{before} truncation. The original
free class-two case is prior; the proposed contribution is the general
finite recursion. Its full word-input implementation exists, with
positive and negative native GAP reconstructions, but the all-rank
structural proof remains a major review priority.
\evidence{N8}{general-proof.md}

\subsection{One fixed group with an undecidable retract problem (N9)}\label{sec:n9}
\entry{N9}{A.~Miasnikov}
Part (a) asks whether the retract problem is decidable in every finitely
generated nilpotent group. Its background distinguishes one fixed
ambient group from the previously established class-wide undecidability
\cite{romankov2016}. Part (b), for free nilpotent groups, has a prior
positive answer~\cite{retracts2017}.

\begin{theorem}\label{prop:n9}
There is a fixed finitely generated torsion-free class-two group $G$
with undecidable retract problem, already on a computable family
\[
 H_n=\langle x_0,x_1^{3n+1}x_2^{-1}\rangle\quad(n\in\N)
\]
of isolated rank-two free class-two subgroups. Their images in both
lower-central layers are primitive sublattices.
\end{theorem}
\begin{proof}
Choose a nonrecursive computably enumerable set $S\subset\N$ containing
$0,1$. DPRM gives one fixed integer polynomial $P$ such that
\[
 n\in S\quad\Longleftrightarrow\quad
 \exists y_1,\ldots,y_m\in\Z\ P(n,y)=0.
\]
The integer version follows by replacing natural-number existential
variables by sums of four squares~\cite{dprm-coq}. Fix an arithmetic
circuit using additions, multiplications and integer constants. Let
$p,o$ be its parameter and output labels.

Use indices $0,1,2,3$ and two new indices $U_z,V_z$ for each scalar
label $z$. Let $L$ be the full integer lattice of alternating matrices
$M$ on these indices satisfying the following \emph{linear} equations:
\begin{align*}
 M_{0,U_z}=M_{1,V_z}&=0,&M_{1,U_z}+M_{0,V_z}&=0,\\
 M_{0,V_z}&=cM_{01}&&\text{for a constant }z=c,\\
 M_{0,V_z}&=M_{0,V_x}+M_{0,V_y}&&\text{for }z=x+y,\\
 M_{U_x,V_y}&=M_{0,V_z}&&\text{for }z=xy,\\
 M_{0,V_o}&=0,\quad M_{03}=M_{12}=0,&
 M_{13}=M_{02}&=3M_{0,V_p}.
\end{align*}
Multiplication is encoded by imposing rank two on $M$, not by nonlinear
equations defining the lattice.

We claim that a rank-two $M\in L$ satisfying
\begin{equation}\label{eq:n9-unit}
 (3n+1)M_{01}-M_{02}=1
\end{equation}
exists precisely when $P(n,y)=0$ is integrally solvable. Put
$d=M_{01}$, $x=M_{02}$ and $t=3n+1$. The Pfaffian on $0,1,2,3$
gives $dM_{23}=x^2$, while $td-x=1$. Hence
\[
 d(M_{23}-t^2d+2t)=1.
\]
Thus $d=\pm1$; since $3\mid x$, the negative sign is impossible.
So $d=1$, $x=3n$ and $M_{0,V_p}=n$. Define
$u_i=M_{i1}$, $v_i=M_{0i}$. All rank-two Pfaffians now give
\[
 M_{ij}=u_iv_j-v_iu_j.
\]
The two coordinates for each scalar are $(u_{U_z},v_{U_z})=(z,0)$
and $(u_{V_z},v_{V_z})=(0,z)$. Thus the remaining linear equations
enforce the circuit over $\Z$. Conversely a circuit solution gives
these coordinate pairs and
\[
 (u_0,v_0)=(1,0),\ (u_1,v_1)=(0,1),\
 (u_2,v_2)=(0,3n),\ (u_3,v_3)=(-3n,0),
\]
which produce the required matrix.

Compute an integer basis $B_1,\ldots,B_s$ of $L$ and define $G$ by
\begin{equation}\label{eq:n9-group}
 [x_i,x_j]=\prod_{\ell=1}^s z_\ell^{B_\ell[i,j]},
 \qquad[x_i,z_\ell]=[z_\ell,z_k]=1.
\end{equation}
On $\Z^D\times\Z^s$ its multiplication is
\[
 (a,c)(b,e)=\left(a+b,c+e-\sum_{i<j}a_jb_iB[i,j]\right).
\]
Bilinearity proves associativity and the normal form realizes the
presentation exactly, proving torsion-freeness. All data are fixed
independently of $n$. Put $a=x_0$, $b_n=x_1^{3n+1}x_2^{-1}$ and
$c_n=[a,b_n]$, with central coordinate
$q_n=(3n+1)B[0,1]-B[0,2]$. A solution at $0\in S$ proves
$q_n\ne0$ for every $n$. Thus $H_n$ is the integral Heisenberg
group with central generator $c_n$.

A retraction $\rho:G\to H_n$ must send $z_\ell$ to
$c_n^{\lambda_\ell}$ and $x_i$ to $a^{r_i}b_n^{s_i}c_n^{k_i}$.
The commutator relations give
\[
 M=\sum_\ell\lambda_\ell B_\ell,\qquad
 M_{ij}=r_is_j-s_ir_j.
\]
Fixing $c_n$ gives~\eqref{eq:n9-unit}, so $M$ has rank two and
$n\in S$. Conversely a rank-two solution of~\eqref{eq:n9-unit}
gives a homomorphism by $x_i\mapsto a^{u_i}b_n^{v_i}$,
$z_\ell\mapsto c_n^{\lambda_\ell}$. Its displayed coordinates fix
$a$ and $b_n$, so it is a retraction. Therefore
$H_n$ is a retract if and only if $n\in S$.

For the asserted isolation, $L$ is the kernel of an integer matrix,
hence saturated. Its basis matrix has an integer right inverse, so
the commutator columns generate $\Z^s=G'$. Matrices $M_0,M_1$ from
solutions at $0,1$ give coefficient vectors $\mu_0,\mu_1$ with
\[
 ((1-n)\mu_0+n\mu_1)\mathbin{\cdot}q_n
 =(1-n)(3n+1)+n(3n-2)=1.
\]
Thus $q_n$ is primitive even when $n\notin S$. The abelianized
subgroup has basis $e_0,(3n+1)e_1-e_2$, a direct summand. Also
$H_n\cap G'=H_n'=\langle c_n\rangle$. If $g^k\in H_n$,
first use the primitive abelianized image to write $g=hz$ with
$h\in H_n$, $z\in G'$, then use primitivity of $q_n$ to conclude
$z\in\langle c_n\rangle$. Hence $g\in H_n$.
\end{proof}

This is an effective construction from a fixed DPRM polynomial, whose
full numerical presentation was not expanded. The implemented small
circuit checks and Lean normalization lemma do not verify DPRM or the
whole reduction. Primitivity of the graded images alone does not make
the subgroup a retract: it does not remove the ambient defining relations.
\evidence{N9}{isolated-inputs-and-prior-scope.md}
